\section{Geometric invariants based on projections}


\paragraph*{Main notations and concepts.}
Let $\R^d$ be the Euclidean $d$-dimensional space and $\Z^d$ its
subset of points with integer coordinates. For any integer $m$, $0 \Le
m \Le d$, we denote by $\mathcal{I}_m$ the set of subsets of
$\{1,\dots,d\}$ with cardinal $m$. For any $J \in \mathcal{I}_m$, the
map $\pi_J: \R^d \to \R^m$ is the orthogonal projection onto the
subspace with coordinates in $J$. Let $p$ be any point of $R^d$ and
let $r$ be a positive value which defines the radius of the local
analysis. We write $B(p,r)$ the $d$-dimensional ball of center $p$ and
radius $r$. The unit $m$-dimensional ball centered on $0$ is denoted
by $B^m$, with volume $\omega_m$. If $T$ is subset of $\R^d$, we write
$B_T(p,r)$ for $B(p,r) \cap T$. The following measures associate a
local value to a set through projections onto sub-axes.

\begin{definition}[local measures]
  \label{def-local-measures}
  Let $\HH^m$ be the $m$-dimensional
  Hausdorff measure. For any subset $Y \subset \R^d$, $p \in Y$, $r>0$, we define the local measures
  \begin{align}
    V_J(Y,p,r) &:= \HH^m\!\left( \pi_J\bigl(B_Y(p,r)\bigr) \right),
    \quad
    && & W_m(Y,p,r) & := \sum_{J \in \mathcal{I}_m} V_J(Y,p,r)^2.
  \end{align}
\end{definition}

We will establish several stability results for these measures:
\begin{enumerate}
\item For $T$ an arbirary $k$-dimensional plane, $W_m(T,p,r)$ is $0$ if
  $m>k$, otherwise is a constant times the square of the volume of the
  $m$-ball, independently of the position and the orientation of $T$ in $\R^d$.
\item For $X$ a smooth closed $k$-manifold, $W_m(X,p,r)$ is close to
  $W_m(T_p X,p,r)$ (for $T_p X$ the tangent plane to $X$ at $p$),
  hence also independent of the orientation of the tangent plane.
\item The same result hold for a covering discretization of $X$.
\item The quantities $V_J$ are also stable and related to the local
  geometry of $X$ or its discretization (tangent vector, normal).
\end{enumerate}
Clearly, these properties give mathematical tools to analyse the local
dimension of arbitrary lattice sets as well as some geometric
information (see section on experiments). Being measures, they are
robust to discretization noise. We will also show the convergence of
these local measures for finer and finer digitizations.

We need a few more definitions. Let $X \subset \R^d$ be a $C^2$-smooth
submanifold of $\R^d$ of dimension $k$ without boundary. We further
suppose it has positive reach $\reach(X)\geq \rho>0$. For $h>0$ a
discretization step, let $Q_h:= \left[-\frac{h}{2},
  \frac{h}{2}\right]^d$. The \emph{$h$-covering discretization of $X$}
is $G_h(X) := \{ z \in h\Z^d~|~X \cap (z+Q_h) \neq \emptyset \}$, its
\emph{$h$-voxelization} is $U_h(X) := \cup_{z \in G_h(X)} (z + Q_h)$.
%%Finally we call \emph{expected $r$-local measure of a $k$-manifold}
%%the quantity $W_m^k := \omega_m^2\binom{k}{m}r^{2m}$.

%% We have immediately:
%% \[
%%     X\subset U_h(X)
%%     \subset
%%     X+\frac{\sqrt d}{2}h B^d.
%% \]

Our main result below states that the local measure of a digitized set
depends essentially on the dimension of its underlying submanifold,
with an error that tends to zero with the discretization
gridstep. This provides a tool for analyzing the local dimension of
digital data.
\begin{theorem}[Local approximation of $W_m$]
\label{thm:main}
For $p \in X$, $0 < m \Le k$, $\forall \alpha \in (0,1)$, there is a constant $C=C(d,k,m,\alpha)$ such that, if $r \Le \alpha\rho$ and $\frac{\sqrt d}{2}h \Le \alpha r$, then
\begin{align*}
  \left|W_m(U_h(X),p,r) - \omega_m^2\binom{k}{m}r^{2m} \right|
  \Le  C \left( h\,r^{2m-1} + \frac{r^{2m+1}}{\rho} \right).
\end{align*}
\end{theorem}
The proof relies on three lemmas: the first one controls the error
coming from the discretization, the second one relates the local
measures on $X$ to the same local measures on the tangent plane, the
third one shows that this value is independent of the orientation of
the tangent plane.

\subsection{Local measures of $k$-balls are invariants}

The fact that the Pythagorean theorem extends to arbitrary volumes has
been rediscovered under different forms several times. The recent
paper of Ancel \cite{ancel2023pythagorean}) provides a general form
for a Lebesgue measurable set $B$ in a $k$-dimensional affine
space. Then $\HH^k(B)^2=\sum_{J \in \mathcal{I}_k} \HH^k(\pi_J(B))^2$ (here
Lebesgue and $k$-th Hausdorff measures are the same).

We provide here a much simpler proof for the case of the ball, which
is sufficient in our context. This formulation also extends the result to lower dimensional projections.
\begin{theorem}\label{thm:ball-invariants}
Let $B^k$ be the unit Euclidean ball embedded isometrically into
$\mathbb{R}^d$ with the linear map $A\in\mathbb{R}^{n\times k}, A^T
A=\mathrm{Id}_k$ (so $k \le d$).  For every integer $m$, $0\le m\le k$, one has
\begin{equation}
  \label{eq:ball-invariants}
  \sum_{J \in \mathcal{I}_m} \HH^m(\pi_J(B^k))^2 = \omega_m^2 \binom{k}{m}.
\end{equation}
The same sum is obviously null if $k < m \Le d$.
%% \end{minipage}
%% \hfill
%% \begin{minipage}{0.45\textwidth}
%% \begin{equation}
%%   \label{eq:pythagorean-ball-2}
%%   W_m(A(rB^k),0,r) = \omega_m^2 \binom{k}{m} r^{2m}.
%% \end{equation}
%% \end{minipage}
\end{theorem}
\begin{proof}
For every subset $J \in \mathcal{I}_m$, Lemma~\ref{lem:svd} (see appendix) applied
to $L=\pi_J\circ A$ gives
\begin{equation}
  \label{eq-vol-proj-ball-1}
  \HH^m(\pi_J( B^k ))^2
  =\omega_m^2 \det(A_J A_J^T),
\end{equation}
where $A_J$ is the matrix containing the rows of $A$ indexed by
$J$. We denote by $\mathcal{I}^k_m$ the set of subsets of
$\{1,\dots,k\}$ with size $m$ (remember that $A_J$ may not be a square
matrix).  By the Cauchy--Binet formula, we have
\begin{equation}
  \label{eq-cauchy-binet-1}
\det(A_J A_J^T)
=
\sum_{I \in \mathcal{I}^k_m}
\det(A_J^I) \det( (A_J^T)_I ),
=
\sum_{I \in \mathcal{I}^k_m}
\det(A_J^I)^2,
\end{equation}
where $A_J^I$ denotes the $m\times m$ submatrix obtained by selecting
the rows indexed by $J$ and the columns indexed by $I$, and where we
have used that $(A_J^T)_I = (A_J^I)^T$ and the equality of the
determinant of a matrix and its transpose.
Therefore, \Equ{eq-vol-proj-ball-1} and \Equ{eq-cauchy-binet-1} implies that
\begin{align}
  \sum_{J \in \mathcal{I}_m}
  \HH^m(\pi_J(B^k))^2
  & =
  \omega_m^2
  \sum_{J \in \mathcal{I}_m}
  \sum_{I \in \mathcal{I}^k_m}
  \det(A_J^I)^2
  \nonumber \\
  \label{eq-proj-vol-ball-2}
  & =
  \omega_m^2
  \sum_{I \in \mathcal{I}^k_m}
  \sum_{J \in \mathcal{I}_m}
  \det(A_J^I)^2.
  && \text{(swapping the sums)}
\end{align}
Fix a subset $I$. The matrix $A^I$ consisting of the corresponding $m$
columns of $A$ satisfies $(A^I)^T A^I=\mathrm{Id}_m$, because the
columns of $A$ are orthonormal.
Applying Cauchy--Binet once more, we get
\begin{equation}
  \label{eq-cauchy-binet-2}
  1 = \det(\mathrm{Id}_m) = \det((A^I)^T A^I) = \sum_{J \in \mathcal{I}_m} \det(A_J^I)^2.
\end{equation}
Since there are $\binom{k}{m}$ possible subsets $I$ in $\mathcal{I}^k_m$, inserting \Equ{eq-cauchy-binet-2} into \Equ{eq-proj-vol-ball-2} gives

\[
\sum_{J \in \mathcal{I}_m}
\HH^m(\pi_J(B^k))^2
=
\omega_m^2
\binom{k}{m}.
\]

This completes the proof. \qed
\end{proof}
Thus the result is independent of the orientation of the ball in the
embedding space, as well as the dimension of this space. Let us denote
$\bar{W}_{m}(k)$ these geometric invariants of the $k$-ball. So for $0
\Le m \Le k$, $\bar{W}_{m}(k)=\omega_m^2\binom{k}{m}$ and for $k < m
\Le d$, $\bar{W}_{m}(k)=0$.  An immediate corollary is that, for $T$
any affine $k$-dimensional plane and $p \in T$, $W_m( B_T(p,r),p,r) =
\bar{W}_{m}(k) r^{2m}$.  These values (scaled by $r^{2m}$) correspond
to the expected values of the local measures when analyzing
digitizations of $k$-dimensional submanifolds (see appendix for tables
expliciting their values in low dimension).





\subsection{Local measures are stable in a neighborhood of a submanifold}

For a positive value $\delta$, we use the notation $(Y)_\delta$ to
denote the dilation of $Y$ with the ball of radius $\delta$. The tangent plane $T_p X$ of $X$ at $p$ is simply written $T$.
\begin{theorem}%%[Stabilité des volumes projetés au voisinage d'une sous-variété]
\label{thm:stabilite-projections}
Under the hypotheses of \RefTheorem{thm:main}, there exists a constant
$\alpha_0 \in (0,1)$ and a constant $C=C(d,k)>0$ such that, for $0<r
\Le \alpha_0 \rho$ and $0 \Le \varepsilon \Le r$, for all
compact set $Y\subset\mathbb R^d$ with
\begin{equation}
\label{eq:encadrement-Y}
    X \cap B(p,r) \subset Y \subset (X \cap B(p,r))_\varepsilon,
\end{equation}
then for any $J\in \mathcal{I}_m$ it holds:
\begin{align}
\label{eq:stabilite-VJ}
\left| V_J(Y,p,r)- V_J(B_T(p,r),p,r)\right|
& \Le C r^{m-1}\left( \varepsilon+\frac{r^2}{\rho} \right), \\
\label{eq:stabilite-W} 
\left| W_m(Y,p,r)-\omega_m^2\binom{k}{m}r^{2m} \right|
& \Le C r^{2m-1} \left( \varepsilon+\frac{r^2}{\rho} \right).
\end{align}
\end{theorem}
\begin{proof}
  To simplify expressions, $C$ is a positive constant depending only
  on $d$ and $k$, whose value may be updated during the proof. We
  write $D_r$ for the disk $B_T(p,r)$.
  \begin{enumerate}
    % \textbf{1. Approximation de $X_r$ par le disque tangent.}
  \item We start by approximating $X_r:=X \cap B(p,r)$ by $D_r$. Since
    $X$ has positive reach, for a sufficiently small $\alpha_0$, the
    set $X$ is in a neighborhood of $p$ the graph of a $C^2$-map
    $f:\Omega\subset T\longrightarrow T^\perp$, with $f(0)=0$,
    $Df(0)=0$, $|f(u)|\Le C\frac{|u|^2}{\rho}$, $\|Df(u)\|\Le
    C\frac{|u|}{\rho}$.
    
    Let $\Phi(u):=p+u+f(u)$. For $x=\Phi(u)\in X_r$, since $u\perp
    f(u)$ we have $|u|\Le |x-p| \Le r$. Furthermore
    $|x-(p+u)|=|f(u)| \Le C\frac{r^2}{\rho}=:\delta$. We obtain:
    \begin{equation}
      \label{eq:X-outer}
      X_r\subset (D_r)_\delta.
    \end{equation}
    Reciprocally, we can set $r_-:=r-C\frac{r^3}{\rho^2}>0$, with
    $\alpha_0$ small enough so that $r-r_- \le \delta$, such that
    \begin{equation}
      \label{eq:inner-graph}
      \Phi(B_T(0,r_-))\subset X_r.
    \end{equation}
    
    Indeed, for $|u|\Le r_-$, $|\Phi(u)-p|^2 = |u|^2+|f(u)|^2 \Le
    |u|^2+C\frac{|u|^4}{\rho^2}$. the choice of $r_-$ induces that
    the last member is smaller or equal to $r^2$.
    
    %\textbf{2. Projection et inclusion extérieure.}
  \item We verify now than the projection of $Y$,
    $\mathcal{F}_J:=\pi_J(Y)$, is included in some dilation of the
    projected disk, denoted by $\mathcal{E}_J:=\pi_J(D_r)$.

    Using \eqref{eq:encadrement-Y} and the fact that
    $\pi_J$ is $1$-Lipschitz, we get $\pi_J(X_r) \subset \mathcal{F}_J
    \subset \pi_J(X_r)+B^m_\varepsilon$.
    Using \eqref{eq:X-outer}, it follows that
    \begin{align}
      \label{eq:outer-projection}
      \mathcal{F}_J\subset \mathcal{E}_J+B^m_\eta,
      \quad\text{with}\quad
      %\label{eq:eta-definition}
      \eta:=\varepsilon+\delta
      =
      \varepsilon+C\frac{r^2}{\rho}.
    \end{align}

    The difficulty lies now in controlling the volume of
    $\mathcal{F}_J$, which must also contain a big enough part of
    $\mathcal{E}_J$. 

    %% \textbf{3. Géométrie de l'ellipsoïde tangent.}
    \item To do so, we must first understand the geometry of
      $\mathcal{E}_J$, which may be degenerated.  Consider the linear
      map $A_J:=\pi_J|_T:T\longrightarrow \R^m$. Let $1 \Ge \sigma_1
      \Ge \cdots \Ge \sigma_m \Ge 0$ the $m$ biggest singular values
      of $A_J$. It is well known that the semi-axes of $\mathcal{E}_J$ are
      $r\sigma_1,\ldots,r\sigma_m$. So its $m$-volume is
      \begin{equation}
        \label{eq:ellipsoid-volume}
        \HH^m(\mathcal{E}_J)
        =
        \omega_m r^m\prod_{i=1}^m\sigma_i.
      \end{equation}
      We separate two cases according to the value of the smallest
      singular value $\sigma_m$.

      %\textbf{4. Cas presque dégénéré.}
    \item If $r\sigma_m\Le C_0\eta$ for $C_0>0$ an arbitrary fixed
      constant, then $A_J$ is almost degenerated and we must show that
      sets $\mathcal{E}_J$ and $\mathcal{F}_J$ are both sufficiently
      small.  In the ellipsoid principal coordinates, $\mathcal{E}_J
      \subset \prod_{i=1}^m [-r\sigma_i,r\sigma_i]$. The $m-1$ first
      dimensions are no greater than $2r$ while the last is at most
      $2C_0\eta$. According to \eqref{eq:ellipsoid-volume}, its volume is bounded as
      \begin{equation}
        \label{eq:degenerate-ellipsoid-volume}
        \HH^m(\mathcal{E}_J)
        \Le C r^{m-1}\eta.
      \end{equation}
      Besides, using \eqref{eq:outer-projection},
      $\mathcal{F}_J\subset\mathcal{E}_J+B^m_\eta$ holds. The right
      set is obviously contained in the box $\prod_{i=1}^m
      [-r\sigma_i-\eta,r\sigma_i+\eta]$.

      Exploiting hypotheses, we see that $\eta \Le C r$. The $m-1$
      first dimensions are of size $O(r)$ while the last one has size
      $O(\eta)$, leading to
      \begin{equation}
        \label{eq:degenerate-F-volume}
        \HH^m(\mathcal{F}_J)
        \Le C r^{m-1}\eta.
      \end{equation}

      Combining \eqref{eq:degenerate-ellipsoid-volume} and
      \eqref{eq:degenerate-F-volume} shows these two sets are small, hence
      close in measure:
      \begin{equation}
        \label{eq:degenerate-result}
        \left|
        \HH^m(\mathcal{F}_J)
        -
        \HH^m(\mathcal{E}_J)
        \right|
        \Le C r^{m-1}\eta.
      \end{equation}
      
      %%\textbf{5. Cas non dégénéré : inclusion intérieure.}
    \item If now $r\sigma_m>C_0\eta$, with $C_0$ big enough, the
      linear map $A_J$ is not degenerate and we will show that the
      projected ellipsoid has a big enough interior.  Let us denote
      $L_J:=(\ker A_J)^\perp \cap T$. We have $\dim L_J=m$ and
      $\widetilde A_J:=A_J|_{L_J}$ is an isomorphism from $L_J$ onto
      $\R^m$. For any radius $s>0$, we have $A_J(B_T(0,s))= \widetilde
      A_J(B_{L_J}(0,s))$.

      Our objective is to understand how the local deformation of $X$
      with respect to $T$ induces a deformation of the projection onto
      $T$. We consider the local map $G_J: B_{L_J}(0,r_-)\to \R^m$
      defined as $G_J(u) := \pi_J(\Phi(u)) = \pi_J(p) + \widetilde
      A_J(u) + \pi_J(f(u))$. Using \eqref{eq:inner-graph},
      \begin{equation}
        \label{eq:G-image}
        G_J(B_{L_J}(0,r_-)) = \pi_J(\Phi(B_{L_J}(0,r_-))) \subset
        \pi_J(X_r) \subset \mathcal{F}_J.
      \end{equation}

      Let us introduce the homotopy $G_{J,t}(u) := \pi_J(p) +
      \widetilde A_J (u) + t\,\pi_J(f(u))$ for $t\in[0,1]$. One can
      see that $G_{J,0}(u) = \pi_J(p)+\widetilde A_J(u)$ and
      $G_{J,1}=G_J$. We use this homotopy to connect the ellipsoid
      $\mathcal{E}_J$ to the projection $\mathcal{F}_J$, both
      slightly retracted. Let us define this slightly retracted
      ellipsoid $\mathcal{E}_J^- := \pi_J(p) + \widetilde
      A_J(B_{L_J}(0,r_-))$.

      From its definition, for any $u$ on the boundary $\partial
      B_{L_J}(0,r_-)$ and for any $t\in[0,1]$, the relation
      $|G_{J,t}(u)-G_{J,0}(u)| \Le |f(u)| \Le \delta$ holds. It induces:
      \begin{equation}
        \label{eq:homotopy-boundary}
        G_{J,t}(\partial B_{L_J}(0,r_-))
        \subset
        (\partial\mathcal{E}_J^-)_{\delta}.
      \end{equation}

      If we choose a point $y$ in the ellipsoid $\mathcal{E}_J^-$, but
      far enough from its boundary, i.e. $\operatorname{dist}
      (y,\partial\mathcal{E}_J^-) >\delta$. Then
      \eqref{eq:homotopy-boundary} implies that the point $y$ cannot
      belong to the image of the boundary $G_{J,t}(\partial
      B_{L_J}(0,r_-))$ for any $t\in[0,1]$.  The Brouwer degree
      $\deg\bigl( G_{J,t}, B_{L_J}(0,r_-), y \bigr)$ is thus well
      defined, constant and independent of $t$ (homotopy invariance of
      degree of an interior point). For $t=0$, the map $G_{J,0}$ is an
      affine isomorphism, its degree is $\pm1$; so the map $G_{J,1}$
      has the same degree, and exactly one antecedent in this
      domain. In particular point $y$ belongs to the image of $G_J$ of
      the domain $B_{L_J}(0,r_-)$. Using \eqref{eq:G-image}, we get
      $y\in \mathcal{F}_J$. Since $y$ was at distance at least
      $\delta$ from the boundary of the contracted ellipsoid, we have
      shown
      \begin{equation}
        \label{eq:inner-contracted}
        \mathcal{E}_J^- \setminus (\partial\mathcal{E}_J^-)_{\delta} \subset \mathcal{F}_J.
      \end{equation}
      It remains to replace $\mathcal{E}_J^-$ with
      $\mathcal{E}_J$. These two ellipsoid are homothetic and, since
      $\pi_J$ is 1-Lipschitz, $d_H(\mathcal{E}_J^-,\mathcal{E}_J) \Le
      r-r_- \Le C\frac{r^3}{\rho^2} \Le C\delta$.  By convexity any
      point of $\mathcal{E}_J$ at distance greater than
      $\delta+(r-r_-)$ from its boundary belongs to $\mathcal{E}_J^-$
      while being at distance greater than $\delta$ from $\partial
      \mathcal{E}_J^-$. Noticing than $\delta + (r-r_-)\Le 2\delta \Le 2\eta$, the
      inclusion \eqref{eq:inner-contracted} shows that
      \begin{equation}
        \label{eq:inner-final}
        \mathcal{E}_J
        \setminus
        (\partial\mathcal{E}_J)_{2\eta}
        \subset \mathcal{F}_J.
      \end{equation}

      %\textbf{6. Contrôle de la couche de bord.}
      \item The exterior inclusion \eqref{eq:outer-projection} and the
        inner inclusion \eqref{eq:inner-final} implies that the
        symmetric difference of $\mathcal{F}_J$ and $\mathcal{E}_J$ is
        $2\eta$-close to the ellipsoid boundary:
        \begin{equation}
          \label{eq:symmetric-difference}
          \mathcal{F}_J\triangle\mathcal{E}_J
          \subset
          (\partial\mathcal{E}_J)_{2\eta}.
        \end{equation}
        It remains to evaluate the $m$-volume of the latter set.  The
        Steiner formula (or Minkowski-Steiner formula) for convex sets
        tells, for $0<t \Le r$ and a convex set $K$ included in a
        $m$-dimensional ball with radius $r$: $\HH^m(K+tB^m) -
        \HH^m(K) \Le C_m t r^{m-1}$. For the erosion $K-tB^m$, the
        argument is that the area of the parallel sets are
        decreasing. Thus $\HH^m(K) - \HH^m (K-tB^m) = \int_0^t
        \mathrm{Area}(K-sB^m) ds \Le \int_0^t \mathrm{Area}(K) ds \Le
        t \mathrm{Area}(K) = m \omega_m tr^{m-1}$. It follows that
        \begin{equation}
          \label{eq:convex-boundary-layer}
          \HH^m((\partial K)_t)
          \Le C'_m t r^{m-1}.
        \end{equation}
        Since $\mathcal{E}_J \subset B^m(\pi_J(p),r)$, applying the
        estimate \eqref{eq:convex-boundary-layer} to $t=2\eta$ gives
        $\HH^m \bigl( (\partial\mathcal{E}_J)_{2\eta} \bigr) \Le C
        \eta r^{m-1}$.  Since $m$-volume are increasing,
        \eqref{eq:symmetric-difference} implies
        \begin{equation}
          \label{eq:nondegenerate-result}
          \left| \HH^m(\mathcal{F}_J) - \HH^m(\mathcal{E}_J) \right|
          \Le C\eta r^{m-1}.
        \end{equation}

        %\textbf{7. Conclusion pour les volumes projetés.}
      \item We can now conclude for both the degenerate case with
        \eqref{eq:degenerate-result} and the non-degenerate case
        with \eqref{eq:nondegenerate-result}, which imply in all cases that
        $\left|\HH^m(\mathcal{F}_J)-\HH^m(\mathcal{E}_J)\right|\Le
        C\eta r^{m-1}$.

        Since $\eta=\varepsilon+C\frac{r^2}{\rho}$ and using the
        definitions of $\mathcal{F}_J$ and $\mathcal{E}_J$, we get:
        \begin{equation*}
          \left| \HH^m(\mathcal{F}_J) - \HH^m(\mathcal{E}_J) \right| =
          \left| V_J(Y,p,r)- V_J(B_T(p,r),p,r)\right| \Le C r^{m-1}
          \left( \varepsilon+\frac{r^2}{\rho} \right),
        \end{equation*}
        which is exactly \eqref{eq:stabilite-VJ}.

        %%\textbf{8. Passage à $W_m^k$.}
      \item It remains to show the relation on $W_m$. Let
        $a_J:=\HH^m(\pi_J(Y))$, $b_J:=\HH^m(\pi_J(D_r))$. The previous
        relation gives $|a_J-b_J| \Le C r^{m-1} \left(
        \varepsilon+\frac{r^2}{\rho} \right)$.  Besides, by hypothesis
        we have $Y\subset (X_r)_\varepsilon$ which is itself included
        in the ball $\subset B(p,r+\varepsilon)$ (a very rough
        estimate). Since $\varepsilon\Le r$, we have $a_J+b_J\Le C
        r^m$. This implies:
        \begin{equation*}
          |a_J^2-b_J^2| = |a_J-b_J|\,(a_J+b_J) \Le C r^{2m-1}
          \left( \varepsilon+\frac{r^2}{\rho} \right).
        \end{equation*}
        Summing on all possible $J \in \mathcal{I}_m$ gives the relation for $W_m$:
        \[
        \left| W_m(Y,p,r)-W_m(D_r,p,r)\right| \Le C r^{2m-1} \left( \varepsilon+\frac{r^2}{\rho}\right).
        \]
        Since $D_r=B(p,r) \cap T$, we have $D_r=p +
        rB^k$. \RefTheorem{thm:ball-invariants} gives:
        \[
        W_m(D_r,p,r) = r^{2m} W_m( B^k,0,1 ) =
        r^{2m} \sum_{J \in \mathcal{I}_m} \HH^m(\pi_J(B^k))^2
        = \omega_m^2\binom{k}{m}r^{2m}.
        \]
        Combining the two above relations prove \eqref{eq:stabilite-W}.
  \end{enumerate}
\end{proof}


\subsection{Discretization error}

We show that the $h$-voxelization of $X$ satisfies the conditions of
the previous stability theorem, with an approximation parameter
$\varepsilon_h=O(\sqrt d\,h)$.

\begin{theorem}[Cubical approximation within a ball]
\label{thm:cubical-admissible-set}
Let $X_r:=X\cap B(p,r), Y_h:=U_h(X)\cap B(p,r)$. There exist universal
constants $\alpha_0>0$ and $C>0$ such that, if $0<r \Le \alpha_0\rho$
and $0<h \Le r$, then $X_r \subset Y_h \subset (X_r)_{\varepsilon_h}$,
with $\varepsilon_h=C\sqrt d\,h$. 
\end{theorem}
\begin{proof}
The first inclusion is immediate due to voxelization being the
covering of $X$ by cubes.  We now prove the second inclusion.  Let
$y\in Y_h$. By definition, $y$ belongs to a cube $Q := z + Q_h$, $z
\in h\Z^d$ such that $Q\cap X\neq\varnothing$.  Choose $x\in Q\cap
X$. Since $Q_h$ has diameter $\sqrt d\,h$,
\begin{equation}
\label{eq:y-x-cube}
    |y-x|\Le \sqrt d\,h.
\end{equation}
Moreover, $y\in B(p,r)$ implies $|x-p| \Le |x-y|+|y-p| \Le r+\sqrt
d\,h$.  Thus $x\in X\cap B(p,r+\sqrt d\,h)$.  If $x\in B(p,r)$, then
$x\in X_r$ and $\operatorname{dist}(y,X_r) \Le |y-x| \Le \sqrt
d\,h$.

It remains to consider the case $r<|x-p|\Le r+\sqrt d\,h$ (where $x$
is slightly outside $X_r$). We must find a point $x_*$ in $X_r$ close to
$x$. Since $r\Le\alpha_0\rho$, with $\alpha_0$ sufficiently small, $X$ is,
in a neighborhood of $p$, the graph over the tangent plane $T:=T_pX$
of a $C^2$ map $f:\Omega\subset T \to T^\perp$ (see proof of
\RefTheorem{thm:stabilite-projections}), with $f(0)=0, Df(0)=0,
\|Df(u)\| \Le C_0\frac{|u|}{\rho}$. Let $\Phi(u):=p+u+f(u)$ and
$x=\Phi(u)$ for some $u$.

Consider the radial curve on $X$ $\gamma(t):=\Phi(tu)$, for $0\Le
t\Le 1$.  The function $g(t):=|\gamma(t)-p|$ is continuous, with
$g(0)=0$ and $g(1)=|x-p|>r$.  Hence there exists $t_*\in(0,1)$ such
that $|\gamma(t_*)-p|=r$.  Set $x_*:=\gamma(t_*)\in X\cap\partial
B(p,r) = X_r$.  We claim that
\begin{equation}
\label{eq:x-xstar}
    |x-x_*| \Le C_1\bigl(|x-p|-r\bigr),
\end{equation}
where $C_1$ is universal provided $\alpha_0$ is sufficiently small.

Indeed, differentiating $\gamma$ wrt $t$ gives $\gamma'(t) = u+Df(tu)[u]$,
so that
\[
    |\gamma'(t)| \Le \left(1+C_0\frac{|u|}{\rho}\right)|u| \Le
    (1+C_0\alpha_0)|u|.
\]
On the other hand, since $f(u)\perp T$, $|\Phi(u)-p|^2 =
|u|^2+|f(u)|^2$ and the same identity holds with $u$ replaced by $tu$.
Using the bounds on $f$ and $Df$, and decreasing $\alpha_0$ if
necessary, the radial function $g(t)=|\Phi(tu)-p|$ satisfies
$g'(t)\Ge c\,|u|$ for some universal constant $c>0$ and all
$t\in[t_*,1]$.

Consequently, $|x-p|-r = g(1)-g(t_*) \Ge c|u|(1-t_*)$. Now
\[
|x-x_*| \Le \int_{t_*}^1|\gamma'(t)|\,dt \Le (1+C_0\alpha_0)|u|(1-t_*) \Le \frac{1+C_0\alpha_0}{c}(|x-p|-r),
\]
using the relation above for the last inequality. This proves
\eqref{eq:x-xstar}. 

Since $|x-p|-r\Le\sqrt d\,h$, it implies $|x-x_*| \Le C_1\sqrt
d\,h$.  Combining this with \eqref{eq:y-x-cube} yields
$\operatorname{dist}(y,X_r) \Le |y-x_*| \Le |y-x|+|x-x_*| \Le
(1+C_1)\sqrt d\,h$.  Choosing $C=1+C_1$ and setting $\varepsilon_h :=
C\sqrt d\,h$, we get $y\in(X_r)_{\varepsilon_h}$.  Since $y\in Y_h$
was arbitrary, $Y_h\subset(X_r)_{\varepsilon_h}$ and we conclude.
\end{proof}



%% We have to take care of the fact that the intersection of $U_h(X)$
%% with the local ball may contain points of $X$ slightly away from this
%% ball.

%% \begin{lemma}%%[Erreur de discrétisation]
%% \label{lem:discretization}
%% Under the hypotheses of \RefTheorem{thm:main}, there is a constant $C=C(d,k,m,\alpha)$ such that, for all $J \in \mathcal{I}_m$:
%% \begin{align*}
%%   \left|V_J(U_h(X),p,r)-V_J(X,p,r)\right| & \Le C h r^{m-1},\\
%% %%  && \quad \text{and} \quad &
%%   \left|W_m(U_h(X),p,r)-W_m(X,p,r)\right| & \Le C h r^{2m-1}.
%% \end{align*}
%% \end{lemma}

%% \begin{proof}
%%   Let $\eps_h=\frac{\sqrt d}{2}h$. It is clear that $X\subset
%%   U_h(X)\subset X+\eps_h B^d$ (for $+$ is the Minkowski sum).  These
%%   inclusions hold after intersection with $B(p,r)$ and after
%%   projection, so:
%%   \begin{align*}
%%     \pi_J(X\cap B(p,r))
%%     \subset
%%     \pi_J(U_h(X)\cap B(p,r))
%%     \subset
%%     \pi_J(X\cap B(p,r+\eps_h))
%%     +
%%     \eps_h B^m,
%%   \end{align*}
%% where $B^m$ is the unit $m$-dimensional ball centered on $0$.  Since
%% $r \Le \alpha\rho$ (smaller than the reach), the submanifold $X$ is a
%% $C^1$ graph $f$ over the tangent plane $T=T_p X$ within $B(p,2r)$,
%% with bounded value $\|f\| \Le C_\alpha\frac{r^2}{\rho}$ and bounded
%% differential $\|Df\| \Le C'_\alpha\frac r\rho$.

%% En particulier, les projections
%% \[
%%     \pi_J(X\cap B(p,r+\eps_h))
%% \]
%% possèdent un contenu de Minkowski $(m-1)$ uniformément contrôlé par
%% \[
%%     C_{d,k,m,\alpha}r^{m-1}.
%% \]
%% Le passage de $r$ à $r+\eps_h$ ajoute donc un volume
%% \[
%%     O_{d,k,m,\alpha}
%%     \left(r^{m-1}\eps_h\right),
%% \]
%% et le voisinage $\eps_h B^m$ ajoute un terme de même ordre.

%% Ainsi
%% \[
%%     \left|
%%     V_J(U_h(X),r)-V_J(X,r)
%%     \right|
%%     \Le
%%     C_{d,k,m,\alpha}
%%     r^{m-1}\eps_h
%%     \Le
%%     C_{d,k,m,\alpha}hr^{m-1}.
%% \]

%% Enfin, toutes les quantités $V_J(X,r)$ et $V_J(U_h(X),r)$ sont
%% $O(r^m)$. Donc
%% \[
%% \begin{aligned}
%% \left|
%% V_J(U_h,r)^2-V_J(X,r)^2
%% \right|
%% &\Le
%% \left|V_J(U_h,r)-V_J(X,r)\right|
%% \\
%% &\qquad\qquad\times
%% \left(V_J(U_h,r)+V_J(X,r)\right)
%% \\
%% &\Le
%% C h r^{2m-1}.
%% \end{aligned}
%% \]
%% La sommation sur les
%% $\binom dm$ ensembles $J$ donne le résultat.
%% \end{proof}
